\ifdefined\MyWebBook\else
\documentclass[../textbook.tex]{subfiles}
\begin{document}
\fi
\bookappendix{实验复现与评估记录模板}{app:experiment-reporting}{3}

本附录提供一套可以直接填写的实验记录结构。它适用于训练、模型比较、RAG、智能体和系统性能实验；不同任务可以删去不适用字段，但保留实验对象、版本、数据、协议、逐例结果和决策依据。统计方法见\bookxref{ch:evaluation-statistics}，数据血缘与划分见\bookxref{ch:data-engineering}，资源口径可查\bookappendixref{app:resource-estimation}。

\section{报告首页：先说明要作出的决定}

实验标题应表达被比较的对象和主要条件，例如“固定检索语料与生成预算下，重排器是否提高引用问答成功率”，而不只写“模型效果测试”。首页至少填写表\ref{tab:appendix-report-overview}中的字段。

\begin{longtable}{@{}p{.24\linewidth}p{.67\linewidth}@{}}
\caption{实验报告首页模板}\label{tab:appendix-report-overview}\\
\toprule
字段 & 填写内容\\\midrule
\endfirsthead
\toprule
字段 & 填写内容\\\midrule
\endhead
实验标识 & 唯一 ID、标题、负责人、开始与结束时间、报告版本。\\
决策问题 & 本实验支持哪项选择、发布、回退或继续研究决定。\\
主要假设 & 预期改变的机制、主要指标的方向及成立条件。\\
实验单位 & 独立抽样和分配的单位，例如用户、会话、文档、问题或训练运行。\\
候选方案 & 基线与候选系统的完整名称；一次只写模型名不足以标识系统。\\
目标总体 & 结果希望推广到的用户、任务、语言、时间范围、负载与硬件范围。\\
主要目标量 & 指标、聚合权重、比较方式与置信区间；说明越大或越小是否更好。\\
最低实用效应 & 足以改变决策的最小提升、非劣界或硬约束阈值。\\
预算与停止规则 & 样本量、运行次数、词元、时间、费用、工具调用及提前停止条件。\\
预注册状态 & 哪些分析在观察最终结果前已经确定；事后分析标为探索性。\\
结论摘要 & 主要估计、区间、约束是否满足及最终决定；与详细结果保持一致。\\
\bottomrule
\end{longtable}

主要假设需要能够被结果否定。例如，“在相同生成模型、提示和预算下，加入重排器使目标任务成功率至少提高两个百分点，且 P95 延迟增加不超过 \SI{100}{\milli\second}”同时规定了对象、质量效应和性能约束。若实验只用于测量现状，则把假设改写为待估目标量及精度要求。

\section{制品与版本清单}

系统版本应覆盖从原始输入到最终输出的完整执行路径。每个制品同时记录人类可读版本和机器可核对标识；标签“最新版”无法支持复现。

\begin{longtable}{@{}p{.23\linewidth}p{.29\linewidth}p{.39\linewidth}@{}}
\caption{实验制品清单模板}\label{tab:appendix-artifact-manifest}\\
\toprule
制品类别 & 身份字段 & 需要记录的内容\\\midrule
\endfirsthead
\toprule
制品类别 & 身份字段 & 需要记录的内容\\\midrule
\endhead
代码 & 仓库、提交、工作区状态 & 提交哈希、未提交补丁或归档包哈希、入口命令、配置文件。\\
模型 & 模型或检查点 ID、哈希 & 架构配置、权重版本、分词器、量化与适配器、合并顺序。\\
数据 & 数据集、快照、清单哈希 & 原始来源、获取时间、许可、过滤、去重、划分和样本计数。\\
提示与模板 & 模板版本、内容哈希 & 系统提示、消息格式、少样本示例、输出模式和变量填充规则。\\
检索资产 & 语料快照、索引版本 & 解析器、切分、嵌入模型、索引参数、重排器和访问过滤。\\
工具与环境 & 工具版本、镜像或环境锁 & 可用工具、权限、沙箱、依赖锁文件、驱动及运行时版本。\\
评审器 & 规则、裁判或验收器版本 & 解析、规范化、评分提示、裁判模型、人工指南和冲突处理。\\
硬件 & 设备型号与拓扑 & 设备数、内存、互连、CPU、存储及相关功率或频率设置。\\
输出 & 运行目录或对象 ID & 原始响应、工具轨迹、日志、逐例分数、聚合结果和报告。\\
\bottomrule
\end{longtable}

哈希确认两个文件的字节身份，不能替代来源、质量和适用性说明。敏感数据无法随报告分发时，应保存受控位置、访问条件、保留期限和可公开的统计摘要，使复核者知道缺失的是访问权限而非版本信息。

\section{数据、划分与抽样}

\begin{table}[htbp]
\centering
\caption{数据与划分记录模板}\label{tab:appendix-data-splits}
\begin{tabularx}{\linewidth}{p{.21\linewidth}p{.20\linewidth}X}
\toprule
数据部分 & 样本或独立组数 & 来源、时间范围、用途与隔离规则\\\midrule
训练集 &  & 拟合参数；记录过滤、混合权重、重复与污染检查。\\
开发集 &  & 选择提示、模型、超参数和错误修复；记录被查看的次数与用途。\\
校准集 &  & 选择阈值、概率校准器或评审器规则；不得再充当最终独立测试。\\
最终测试集 &  & 仅评估冻结方案；记录访问控制、冻结时间与任何意外暴露。\\
压力与安全集 &  & 覆盖长尾、边界和高损失失败；报告构成，不与自然分布样本混成一个比例。\\
线上观察或实验流量 &  & 合格条件、随机化单位、曝光窗口、排除规则和样本比例检查。\\
\bottomrule
\end{tabularx}
\end{table}

对每次划分补充以下说明：
\begin{enumerate}
 \item 抽样框与目标总体如何对应，哪些对象没有进入抽样框。
 \item 独立单位是什么；同一用户、文档、模板或来源族是否跨集合出现。
 \item 去重采用哪些键、相似度和阈值，近重复检查是否跨训练与测试执行。
 \item 样本权重、分层或过采样如何进入最终聚合，报告的是样本平均还是目标总体加权平均。
 \item 缺失标签、无法执行、超时和不适用项如何计入分子与分母。
 \item 数据版本变化后，哪些结果仍可比较，哪些需要重新运行。
\end{enumerate}

\section{系统与运行协议}

\subsection{生成与智能体实验}

\begin{longtable}{@{}p{.25\linewidth}p{.66\linewidth}@{}}
\caption{生成系统与智能体运行配置}\label{tab:appendix-generation-protocol}\\
\toprule
字段 & 填写内容\\\midrule
\endfirsthead
\toprule
字段 & 填写内容\\\midrule
\endhead
输入构造 & 消息顺序、模板、截断、上下文上限、附件预处理和无效输入处理。\\
生成条件 & 温度、Top-$p$/Top-$k$、最大输出、停止条件、随机种子和每题尝试次数。\\
选择策略 & 单次输出、最佳分数、自洽投票、验证器选择或人工挑选；选择预算单列。\\
工具协议 & 工具集合、参数模式、授权、调用上限、超时、重试、失败返回和副作用隔离。\\
状态协议 & 会话记忆、缓存、持久化、环境重置和跨案例状态是否清空。\\
检索协议 & 查询构造、候选数、过滤、重排、证据上限、索引快照和无结果处理。\\
安全控制 & 拒绝、回退、人工确认、策略模型和中止后的计分方式。\\
并发与速率 & 客户端并发、服务端批处理、速率限制、请求到达模式和重试退避。\\
失败分类 & 无效输出、超时、工具错误、环境错误、拒答和验收失败的互斥或层级规则。\\
\bottomrule
\end{longtable}

随机种子应与实际随机源一起记录，包括数据顺序、参数初始化、采样和环境随机性。固定种子便于重放同一随机过程；跨设备、框架、内核或并行布局时，浮点求和顺序可能变化，因此还要声明允许的数值或指标容差。

\subsection{训练实验}

训练记录至少包括：
\begin{itemize}
 \item 初始化检查点、分词器、可训练与冻结参数集合、参数总量和可训练参数量；
 \item 数据混合、采样权重、有效词元定义、序列打包和标签屏蔽规则；
 \item 优化器、学习率调度、权重衰减、梯度裁剪、损失缩放与正则项；
 \item 每设备微批、梯度累积、全局批量、设备数和分布式并行配置；
 \item 总更新步、处理样本与有效词元数、验证频率、早停与最佳检查点选择规则；
 \item 精度格式、激活重计算、注意力内核、编译选项和确定性设置；
 \item 检查点内容、保存周期、恢复位置，以及优化器、调度器和随机状态是否恢复；
 \item 训练损失、验证指标、梯度与参数统计、吞吐、峰值显存、故障和恢复事件。
\end{itemize}

全局批量的定义写成可复核公式。例如数据并行组大小 $N_d$、每设备微批 $B_\mu$、梯度累积步数 $G$ 时，按样本计的全局批量为 $N_dB_\mu G$；按有效词元计的批量还要根据各序列非填充位置求和。流水并行的在途微批不自动成为额外的优化器批量。

\subsection{系统性能实验}

性能报告需冻结负载与测量边界：
\begin{itemize}
 \item 请求输入长度、输出长度、并发或到达率的分布，以及预填充和解码是否分别统计；
 \item 模型、数值格式、并行、批处理、缓存、调度和服务副本配置；
 \item 冷启动或预热过程、计时起止点、同步方法、采样窗口和重复次数；
 \item 首词元时间、逐输出词元时间、端到端延迟、吞吐及采用的分位数算法；
 \item 峰值显存、主机内存、设备利用率、通信量、错误率、排队和丢弃；
 \item 功耗或成本的计价边界，包括空闲资源、失败重试及共享基础设施分摊方式。
\end{itemize}

达到固定并发与按固定到达率施压回答的是不同问题。前者观察给定在途请求数，后者还会产生排队和过载行为。对比方案必须采用同一种负载生成方式，并同时报告质量约束，避免通过缩短输出或增加拒绝降低延迟。

\section{指标与统计计划}

在运行最终实验前填写表\ref{tab:appendix-metric-plan}。主要指标用于核心决策，次要指标解释取舍，诊断指标定位机制；三者不应在观察结果后互换角色。

\begin{longtable}{@{}p{.15\linewidth}p{.17\linewidth}p{.17\linewidth}p{.17\linewidth}p{.21\linewidth}@{}}
\caption{指标与统计计划模板}\label{tab:appendix-metric-plan}\\
\toprule
指标 & 逐例定义 & 聚合与权重 & 不确定性方法 & 决策规则\\\midrule
\endfirsthead
\toprule
指标 & 逐例定义 & 聚合与权重 & 不确定性方法 & 决策规则\\\midrule
\endhead
主要质量指标 &  &  &  & 最低实用效应或非劣界\\
高损失失败 &  &  &  & 最大允许率或零容忍规则\\
延迟 & 起止点与单位 & 中位数、P95 等 & 重复或簇重采样 & 上限或相对变化\\
资源或费用 & 包含项 & 每请求、每成功任务等 & 价格与负载情景 & 预算上限\\
次要指标 &  &  &  & 解释性，不替换主要结论\\
切片指标 & 切片规则 & 样本数与权重 & 区间或描述性结果 & 预设或探索性标记\\
\bottomrule
\end{longtable}

统计计划还应写明：
\begin{enumerate}
 \item 独立抽样单位和配对键。相同案例比较两个系统时保留逐例配对；同一用户多条记录按用户或其他独立簇处理。
 \item 区间或检验方法、置信水平、Bootstrap 重复次数、随机种子及分位数约定。
 \item 多次运行和多随机种子的聚合层级；案例变异与运行变异分别对应什么总体。
 \item 多指标、多切片和反复查看结果时采用的控制方法；探索性发现在哪个新样本上复验。
 \item 样本量依据，包括希望检测的效应、基率、方差、功效、聚类和预期损耗。
 \item 无效结果、缺失、超时和提前停止如何进入统计量。
\end{enumerate}

\section{结果表与失败证据}

\subsection{主结果}

\begin{table}[htbp]
\centering
\caption{候选方案比较的主结果模板}\label{tab:appendix-main-results}
\begin{tabularx}{\linewidth}{p{.19\linewidth}p{.15\linewidth}p{.15\linewidth}p{.19\linewidth}X}
\toprule
指标 & 基线 & 候选 & 配对差或比值 & 区间、样本数与判定\\\midrule
主要质量指标 &  &  &  &  \\
高损失失败率 &  &  &  &  \\
P50 延迟 &  &  &  &  \\
P95 延迟 &  &  &  &  \\
每成功任务费用 &  &  &  &  \\
其他硬约束 &  &  &  &  \\
\bottomrule
\end{tabularx}
\end{table}

表中同时给出点估计、区间和有效样本或独立簇数。仅报告“提升 $x\%$”时无法区分百分点差、相对变化和比值，应明确写成例如“从0.60升至0.66，绝对增加0.06，相对增加10\%”。多个方案必须绑定各自完整版本，不以展示名称代替配置。

\subsection{切片、消融与失败分析}

切片结果至少列出切片定义、样本数、总体权重、各方案结果与区间。总体变化可能来自切片内性能变化，也可能来自切片构成变化，两部分应分别检查。事后发现的切片保留用于诊断，并在独立材料上复验后再支持发布结论。

消融实验每次删除或替换一个明确组件，同时固定其余条件和预算。若组件改变了输出长度、工具次数或成功任务数量，应同时报告这些中介量。多个组件存在交互时，逐个单独消融不能还原全部组合效果，需要补充组合设计或明确结论只适用于当前基线。

失败分析应保留逐例证据，并使用预先定义或迭代版本化的分类体系：
\begin{longtable}{@{}p{.15\linewidth}p{.15\linewidth}p{.18\linewidth}p{.17\linewidth}p{.22\linewidth}@{}}
\caption{失败案例记录模板}\label{tab:appendix-failure-log}\\
\toprule
案例 ID & 切片与严重度 & 观察到的失败 & 证据与执行阶段 & 归因状态及后续验证\\\midrule
\endfirsthead
\toprule
案例 ID & 切片与严重度 & 观察到的失败 & 证据与执行阶段 & 归因状态及后续验证\\\midrule
\endhead
 &  &  &  &  \\
 &  &  &  &  \\
 &  &  &  &  \\
\bottomrule
\end{longtable}

“回答错误”是观察结果，不是根因。检索未召回、上下文组装错误、模型忽略证据、工具参数错误、环境执行失败和验收器误判属于不同阶段。根因归因需要通过重放、替换组件或最小反例验证；尚未验证时记录为候选原因。

\section{复现包与交付检查}

一个可复核的实验包建议包含：
\begin{enumerate}
 \item 只读的实验配置和制品清单，包括所有哈希、版本、参数及环境锁；
 \item 从数据清单生成实际输入的脚本或确定性步骤，以及数据访问说明；
 \item 启动基线和候选方案的命令、随机种子列表、预算与停止规则；
 \item 原始逐例输出、工具轨迹、时间戳、错误和资源日志；
 \item 从原始记录计算指标、区间、表格和图形的分析代码；
 \item 本报告、已知偏差、异常运行、排除案例及排除理由；
 \item 在隔离环境中完成一次最小重放的结果与容差判定。
\end{enumerate}

交付前依次确认：数据和最终测试没有参与方案选择；基线与候选使用相同案例和预算；版本与输出能够一一追溯；主要指标、区间和决策规则均已预先定义；失败与超时没有从分母中静默消失；性能测量包含负载和计时边界；结论范围没有超过目标总体和实验条件；报告中的表格可以从保存的逐例记录重新生成。

\begin{note}[最小充分报告]
若篇幅受限，正文至少保留决策问题、系统与数据版本、主要协议、样本或独立簇数、点估计与区间、预算、主要失败及结论范围。完整配置、逐例结果和分析代码放入可追溯的复现包，并从报告链接到对应版本。
\end{note}

\chapterreferences
\myendchapterdocument
